% !PRINTED 780
\weberpage{780}
% running head "II. Nachtrag." at the top of this page is not transcribed as body text.
%<<CONT
\[
Q(s) = \frac{1}{m(s-1)} \overset{A}{\Sigma} \chi(A) + \overset{A}{\Sigma} \chi(A)\, C(s),
\]
und daraus der Satz:

2. \gesperrt{Die Summen $Q_2, \ldots, Q_\mu$ erhalten für $s=1$ endliche
Werthe,} $Q_1$ \gesperrt{wird unendlich, und zwar wird}
\[
\mathrm{Lim}_{s=1} (s-1)\, Q_1(s) = \frac{\mu}{m}.
\]

Um die Summen $Q(s)$ umzuformen, bezeichnen wir mit $p$ jede in $m$ nicht
aufgehende Primzahl und multipliciren alle die unendlichen Reihen von der
Form
\[
1 + \frac{\chi(p)}{p^s} + \frac{\chi(p^2)}{p^{2s}} + \cdots
= \frac{1}{1 - \chi(p)\, p^{-s}}
\]
mit einander. Dadurch erhalten wir [vgl. \S.\,192, (8)]:
\weq{11}{Q(s) = \overset{p}{\Pi} \frac{1}{1 - \chi(p)\, p^{-s}}.}
Wenn nun $p$ zum Exponenten $f$ gehört, wenn also $f$ die kleinste positive,
der Congruenz
\weq{12}{p^f \equiv 1 \pmod{m}}
genügende Zahl, und $\mu = ef$ ist, so kommt unter den $\chi(p)$ jede
$f^{\text{te}}$ Einheitswurzel genau $e$ mal vor, und es ist also (für ein
variables $x$):
\[
\overset{\chi}{\Pi} [1 - \chi(p)\, x] = (1 - x^f)^e,
\]
demnach ergiebt sich aus (11):
\weq{13}{Q_1 Q_2 \ldots Q_\mu = \overset{p}{\Pi} \frac{1}{(1-p^{-fs})^e}.}

Wir führen nun, ohne die Irreducibilität der Kreistheilungsgleichung
vorauszusetzen, den Kreistheilungskörper $\Omega_m$ ein, dessen noch
unbekannten Grad wir mit $\nu$ bezeichnen wollen. Da aber die primitive
$m^{\text{te}}$ Einheitswurzel $r$ einer rationalen Gleichung
$\mu^{\text{ten}}$ Grades genügt, so ist
\weq{14}{\nu \leqq \mu.}

Die Zahlen $\omega$ des Körpers $\Omega_m$ sind alle in der Form
\weq{15}{a\omega = a_0 + a_1 r + \cdots + a_{\nu-1} r^{\nu-1}}
mit ganzen rationalen Coëfficienten $a, a_0, a_1, \ldots, u_{\nu-1}$ % [sic] -- printed $u_{\nu-1}$; display (15) has $a_{\nu-1}$
darstellbar, und wenn $\omega$ eine \gesperrt{ganze Zahl} ist, so können wir
für $a$ eine feste ganze Zahl setzen, deren nähere Kenntniss nicht erforder-
%CONT>>
